Fix an identifying cell \(\psi_j\), put
\[
\mathbf p_{\mathrm{obs}}=(P_{\mathrm{obs}}(o):o\in\mathcal X_{\mathcal O}),
\]
and suppose that \(\psi_j(\mathbf p_{\mathrm{obs}})\) holds.  Fix
\(\mathcal M\in\mathfrak F(\mathcal G,P_{\mathrm{obs}})\).  By
\(\mathrm{def{:}observational\text{-}fiber}\), this model is in
\(\mathfrak M_{+}^{\mathrm{rm}}(\mathcal G)\) and induces
\(P_{\mathrm{obs}}\).  The forward implication of
\(\mathrm{thm{:}response\text{-}fiber\text{-}equivalence}\) therefore supplies a response vector
\(\theta_{\mathcal M}\) satisfying
\[
\mathsf C_{\mathcal G}(\theta_{\mathcal M};\mathbf p_{\mathrm{obs}})
\quad\text{and}\quad
q_{a,y}(\theta_{\mathcal M})=Q_a^{\mathcal M}(y)
\quad
\text{for every }(a,y)\in\mathcal X_A\times\mathcal X_Y.
\tag{1}
\]
In particular, the feasible response fiber over \(\mathbf p_{\mathrm{obs}}\) is nonempty.  The strict inequalities in
\(\mathsf C_{\mathcal G}\) also state that every full-data cell polynomial is positive, as required by
\(\mathrm{ass{:}strict\text{-}positivity}\).

Fix \((a,y)\in\mathcal X_A\times\mathcal X_Y\).  Because \(\psi_j\) is identifying for this coordinate,
\(\mathsf I_{a,y}(\mathcal G;\mathbf p_{\mathrm{obs}})\) is true.  Expanding
\(\mathrm{def{:}terminal\text{-}qe\text{-}test}\), this means
\[
\neg\exists\theta^0\,\exists\theta^1
\left[
\mathsf C_{\mathcal G}(\theta^0;\mathbf p_{\mathrm{obs}})
\wedge
\mathsf C_{\mathcal G}(\theta^1;\mathbf p_{\mathrm{obs}})
\wedge
q_{a,y}(\theta^0)\ne q_{a,y}(\theta^1)
\right].
\tag{2}
\]
Thus any two feasible response vectors have the same \(q_{a,y}\)-value.  By (1), at least one feasible vector exists.  Hence the relation
\[
\mathsf V_{a,y}(\mathbf p_{\mathrm{obs}},z)
\quad\Longleftrightarrow\quad
\exists\theta\left[
\mathsf C_{\mathcal G}(\theta;\mathbf p_{\mathrm{obs}})
\wedge z=q_{a,y}(\theta)
\right]
\tag{3}
\]
has exactly one real solution, denoted \(z_{a,y}\).  Substituting \(\theta_{\mathcal M}\) from (1) into (3) gives
\[
z_{a,y}=q_{a,y}(\theta_{\mathcal M})=Q_a^{\mathcal M}(y).
\tag{4}
\]

It remains to check that either allowed representation of the cell functional evaluates to this unique value.  If the cell stores the graph relation (3), then
\(\mathrm{def{:}symbolic\text{-}observed\text{-}partition}\) and the definition of the successful output give
\[
f_{a,j}(P_{\mathrm{obs}})(y)=z_{a,y}.
\tag{5}
\]
Otherwise, \(\mathrm{def{:}tcf\text{-}procedure}\) stores a rational trace
\(r(\mathbf p)=N(\mathbf p)/D(\mathbf p)\) only after verifying on the whole cell that
\[
\forall\mathbf p\,\forall z\,
\left[
\psi_j(\mathbf p)\wedge\mathsf V_{a,y}(\mathbf p,z)
\Longrightarrow
D(\mathbf p)\ne0\wedge D(\mathbf p)z=N(\mathbf p)
\right].
\tag{6}
\]
Apply (6) to \(\mathbf p=\mathbf p_{\mathrm{obs}}\) and \(z=z_{a,y}\).  The premises hold by the choice of the cell and (3), so
\[
D(\mathbf p_{\mathrm{obs}})\ne0,
\qquad
D(\mathbf p_{\mathrm{obs}})z_{a,y}=N(\mathbf p_{\mathrm{obs}}).
\tag{7}
\]
The first relation in (7) licenses division; the second then yields
\[
r(\mathbf p_{\mathrm{obs}})=z_{a,y}.
\tag{8}
\]
The annotated-cell definition of \(f_{a,j}\), together with (5) in the graph-relation case and (8) in the rational case, proves
\[
f_{a,j}(P_{\mathrm{obs}})(y)
=z_{a,y}
=Q_a^{\mathcal M}(y).
\tag{9}
\]
The choices of \(\mathcal M\), \(a\), and \(y\) were arbitrary, so (9) proves the coordinatewise assertion throughout the entire compatible positive fiber, rather than merely at a response-vector oracle.

We record explicitly the direct score-contrast bridge used by the target-only refinement.  For a feasible response vector define the proof intermediate
\[
t_{s_Y}(\theta)
=\sum_{y\in\mathcal X_Y}s_Y(y)
\left\{q_{a_1,y}(\theta)-q_{a_0,y}(\theta)\right\}.
\tag{10}
\]
The sum is finite because membership in \(\mathfrak M_{+}^{\mathrm{rm}}(\mathcal G)\), as specified by \(\mathrm{def{:}positive\text{-}rm\text{-}class}\), makes \(\mathcal X_Y\) finite.  Applying (1) at \(a_1\) and \(a_0\), and then using \(\mathrm{def{:}causal\text{-}query}\), gives
\[
\begin{aligned}
t_{s_Y}(\theta_{\mathcal M})
&=\sum_{y\in\mathcal X_Y}s_Y(y)
\left\{Q_{a_1}^{\mathcal M}(y)-Q_{a_0}^{\mathcal M}(y)\right\}\\
&=\tau_{s_Y}.
\end{aligned}
\tag{11}
\]
Consequently, if a cell satisfies the direct constancy formula
\[
\neg\exists\theta^0\,\exists\theta^1
\left[
\mathsf C_{\mathcal G}(\theta^0;\mathbf p)
\wedge
\mathsf C_{\mathcal G}(\theta^1;\mathbf p)
\wedge
t_{s_Y}(\theta^0)\ne t_{s_Y}(\theta^1)
\right],
\tag{12}
\]
then at \(\mathbf p_{\mathrm{obs}}\) its unique-value relation
\[
\exists\theta\left[
\mathsf C_{\mathcal G}(\theta;\mathbf p_{\mathrm{obs}})
\wedge z=t_{s_Y}(\theta)
\right]
\tag{13}
\]
has the unique solution \(z=\tau_{s_Y}\), by (1), (11), and (12).  The same argument as (6)--(8) shows that any rational annotation certified by the corresponding whole-cell implication also equals \(\tau_{s_Y}\).  This argument treats the finitely many values \(s_Y(y)\) as supplied real coefficients; it neither assumes that they are rational nor requires a finite encoding of them for the mathematical equality.

Finally, the displayed consequence in the theorem follows directly from (9): finiteness permits termwise substitution, so
\[
\begin{aligned}
&\sum_{y\in\mathcal X_Y}s_Y(y)
\left\{f_{a_1,j}(P_{\mathrm{obs}})(y)-f_{a_0,j}(P_{\mathrm{obs}})(y)\right\}\\
&\qquad=
\sum_{y\in\mathcal X_Y}s_Y(y)
\left\{Q_{a_1}^{\mathcal M}(y)-Q_{a_0}^{\mathcal M}(y)\right\}
=\tau_{s_Y},
\end{aligned}
\tag{14}
\]
where the last equality is \(\mathrm{def{:}causal\text{-}query}\).  All divisions used above are justified by the nonzero-denominator implication (6), and all conditioning and support needed to construct the response fiber is contained in the positive-fiber equivalence (1).  This completes the proof.